\begin{scanpage}{0203}{185}
\noindent a corresponding semi-infinite conductor. It is not difficult to establish this by investigating the radiation of a semi-infinite conductor excited by a concentrated emf.
\sectiontitle{30}{Primary and Secondary Diffraction by a Passive Vibrator}
Let a plane electromagnetic wave fall at an angle $\vartheta_0$ on a thin cylindrical conductor of length $L=z_2-z_1$ and radius $a$ (Figure 73). For purposes of generality, we will consider that the incident wave's electric field $E_0$ forms an angle $\alpha$ with the plane of the figure. Then, its tangential component on the conductor surface will equal
\[E_z^e=E_{0z}\cdot e^{iw_0z},\tag{30.01}\]
where
\[E_{0z}=E\sin\vartheta_0,\quad E=E_0\cos\alpha,\quad w_0=-k\cos\vartheta_0.\tag{30.02}\]
The current induced in the vibrator by this field was investigated by us in \S\ 28. As was already indicated above, Expression (28.14) which was obtained for it has a relatively low precision near the conductor ends. Therefore, it is inadvisable to seek the fringing field by integration of the current. Let us also note that the fringing field found by such a method does not satisfy the principle of duality.

We shall seek the scattering characteristic of a passive vibrator by starting from the following scattering picture which naturally follows from the previous results. An incident plane wave, being diffracted at the conductor ends, excites primary edge waves which are radiated into the surrounding space. Being propagated along the conductor, each of these waves experiences diffraction at the opposite end of the conductor and excites secondary edge waves. The latter, in turn, generate tertiary edge waves, etc. The total fringing field is comprised of the sum of all the edge waves being formed during sequential (multiple) diffraction at the conductor's ends.
\end{scanpage}
